% =====================================================================
% --- CAPITOLO 4: VERIFICA E VALIDAZIONE ---
% =====================================================================
\chapter{Verifica e Validazione}
\label{cap:testing}
% Riferimento Traccia: "Output: Testing e valutazione dell’usabilità"

Questo capitolo descrive le strategie di verifica (statica e dinamica) e validazione adottate per garantire la qualità del software.

\section{Verifica Statica (Analisi del Codice)}
La verifica statica è stata condotta mediante l'utilizzo di \textbf{SonarQube}. Questa attività rientra nelle strategie di \textbf{Fault Avoidance}, mirando a rilevare difetti nel codice sorgente prima dell'esecuzione.

L'analisi automatizzata ha permesso di monitorare le seguenti metriche di qualità:
\begin{itemize}
    \item \textbf{Code Smells:} Identificazione di strutture di codice che, pur non essendo bug, violano i principi di design (es. metodi troppo lunghi, codice duplicato) e riducono la manutenibilità.
    \item \textbf{Cyclomatic Complexity:} Misurazione della complessità del flusso di controllo. SonarQube segnala metodi con complessità ciclomatica elevata, suggerendo refactoring per semplificare i percorsi di esecuzione e facilitare il testing.
    \item \textbf{Reliability \& Security:} Rilevamento preventivo di potenziali bug logici e vulnerabilità di sicurezza (es. injection flaws).
\end{itemize}

\textit{[Placeholder per screenshot SonarQube (Quality Gate, Bugs, Vulnerabilities, Code Smells).]}

\section{Verifica Dinamica (Unit Testing)}
% Riferimento Traccia: "Codice xUnit per unit testing di quattro... metodi non banali..."
% Riferimenti Slide: 27 - 29 - Verification and Validation.pdf

Per la verifica dinamica è stato utilizzato il framework \textbf{Jest}. La strategia adottata è quella del \textbf{Black-Box Testing}: i test sono stati progettati basandosi esclusivamente sulle specifiche funzionali (input e output atteso), ignorando la struttura interna del codice.

\subsection{Metodo 1: [Nome Metodo Es. createAppointment]}
\begin{description}
    \item[Descrizione Metodo] Il metodo \texttt{createAppointment(dto, userId)}...
    \item[Classi di Equivalenza (Input)]
    Per la selezione dei casi di test, è stata applicata la tecnica del partizionamento in \textbf{Classi di Equivalenza}:
    \begin{itemize}
        \item \textbf{dto.date}: [CE1: Data futura valida], [CE2: Data passata (invalida)], [CE3: Data in formato non valido (invalida)]
        \item \textbf{userId}: [CE1: ID utente valido], [CE2: ID utente non esistente (invalido)]
    \end{itemize}
    \item[Oracolo del Test]
    Le asserzioni di Jest (es. \texttt{expect(...).toBe(...)}) fungono da \textbf{Oracolo}, verificando che l'output effettivo del sistema coincida con l'output atteso definito nelle specifiche.
    \item[Criteri di Copertura] È stato adottato il criterio R-WECT (Robust Weak Equivalence Class Testing).
\end{description}
% Inserisci qui lo snippet di codice del test
% \begin{lstlisting}[language=JavaScript, caption={Test per createAppointment}]
% describe('AppointmentsService', () => {
%   it('should not create an appointment for a past date', async () => {
%     // ...
%     await expect(service.createAppointment(pastDto, userId))
%       .rejects.toThrow(BadRequestException);
%   });
%   // ...
% });
% \end{lstlisting}

% ... Replica per gli altri metodi di test ...

\section{Valutazione dell'Usabilità}
% Riferimento Traccia: "Descrizione delle tecniche di valutazione dell’usabilità applicate"

\subsection{Expert Reviews / Inspections}
% Riferimento Traccia: "definizione di una checklist..."
È stata condotta un'ispezione esperta (Expert Review) basata sulle \textbf{8 Regole d'Oro di Shneiderman} (come menzionato nel progetto \texttt{docIng1.pdf, pag. 74}).
\begin{longtable}{|L{12cm}|C{1cm}|C{1cm}|}
\hline
\textbf{Euristica di Shneiderman} & \textbf{Sì} & \textbf{No} \\ \hline
\endfirsthead
\multicolumn{3}{c}%
{{\bfseries \tablename\ \thetable{} -- segue dalla pagina precedente}} \\ \hline
\textbf{Euristica} & \textbf{Sì} & \textbf{No} \\ \hline
\endhead
\hline \multicolumn{3}{r}{{\emph{continua nella pagina seguente}}} \\ \endfoot
\hline
\endlastfoot
1. Sforzarsi per la coerenza & \checkmark & \\ \hline
2. Abilitare scorciatoie per gli utenti esperti & & \checkmark \\ \hline
3. Fornire un feedback informativo & \checkmark & \\ \hline
4. Progettare dialoghi per chiudere un'azione & \checkmark & \\ \hline
5. Prevenire gli errori & \checkmark & \\ \hline
6. Consentire annullamenti semplici delle azioni & & \checkmark \\ \hline
7. Favorire il controllo da parte dell'utente & \checkmark & \\ \hline
8. Ridurre il carico cognitivo & \checkmark & \\
\hline
\end{longtable}
% Discuti qui i risultati (es. "Come emerso, l'app eccelle nel feedback ma manca di scorciatoie...")

\subsection{Esperimento con Utenti Reali}
% Riferimento Traccia: "Progettazione e conduzione di un esperimento con utenti reali..."

\subsubsection{Descrizione dei Soggetti, Procedura e Metriche}
\begin{description}
    \item[Soggetti Reclutati] 5 utenti (3 studenti di informatica, 2 studenti di altre facoltà), età 20-25.
    \item[Procedura Sperimentale] Ad ogni utente è stato chiesto di completare 3 task sull'applicazione: (1) Trovare un bilocale a Napoli, (2) Prenotare una visita, (3) Inviare un'offerta.
    \item[Metriche Considerate] \textit{(Vedi docIng1.pdf, pag. 81)}
    \begin{itemize}
        \item \textbf{Task Completion Rate (Efficacia):} Percentuale di utenti che hanno completato il task.
        \item \textbf{Task Completion Time (Efficienza):} Tempo medio per completare il task.
        \item \textbf{Soddisfazione Utente:} Punteggio medio del survey post-test.
    \end{itemize}
\end{description}

\subsubsection{Definizione del Survey}
% Riferimento Traccia: "Definizione di un survey..."
Agli utenti è stato sottoposto un questionario post-test per valutare la soddisfazione (scala Likert 1-5):
\begin{enumerate}
    \item Quanto hai trovato facile trovare un immobile? (1=Molto Difficile, 5=Molto Facile)
    \item L'interfaccia dell'applicazione è chiara e intuitiva? (1-5)
    \item Il processo di prenotazione visita è stato fluido? (1-5)
    \item Hai trovato le previsioni meteo utili durante la prenotazione? (1-5)
\end{enumerate}

\subsubsection{Presentazione e Discussione dei Risultati}
% Riferimento Traccia: "Presentazione e discussione dei risultati..."
% Inserisci qui grafici e analisi (vedi docIng1.pdf, pag. 84-87)
% \begin{figure}[h!]
%     \centering
%     \includegraphics[width=0.7\textwidth]{images/survey-results.png}
%     \caption{Risultati del Survey di Usabilità.}
% \end{figure}
\textit{[Placeholder per grafici e discussione dei risultati.]}
% Es. "Il 100% degli utenti ha completato i task 1 e 3. Il Task 2 (Prenotazione) ha mostrato un tempo medio di... I risultati del survey (Fig. X) mostrano un'alta soddisfazione (media 4.5/5)..."